\pdfbookmark[2]{Zu VIII.}{nexus-remarks-viii}
\begin{center}\large Zu VIII.\end{center}

Hier bedeutet $p$ eine positive Primzahl von der Form $4n+1$; die Bezeichnung stimmt mit der in der Abhandlung \emph{Theoria residuorum biquadraticorum} I. art. 23 angewendeten überein; es ist also
\[
\begin{gathered}
f\equiv1\pdot2\pdot3\dotsfour\tfrac12(p-3)\pdot\tfrac12(p-1)\pmod p\\
p=aa+bb;\quad a\equiv1\pmod4;\quad b\equiv af\pmod p
\end{gathered}
\]
die mit $\alpha,\beta$ bezeichneten Zahlen sind durch die Zerlegung $p=\alpha\alpha+2\beta\beta$ bestimmt. Die Columne $f$ ist den beiden vorgefundenen Tabellen hinzugefügt; ausserdem sind einige Lücken in denselben ausgefüllt.

Der im Text aufgestellte Satz hängt mit dem biquadratischen Charakter der Zahl $2$ zusammen; da nämlich (vergl. \emph{Theoria resid. biqu.} I. art. 21)
\[
2^{\frac{p-1}4}\equiv f^{\frac12 b}\pmod p
\]
ist, so folgt aus der Congruenz
\[
b\equiv2m+a-1\pmod8
\]
die andere
\SourceLeaf{310}
\[
2^{\frac{p-1}4}\equiv f^{m+\frac{a-1}2}\pmod p
\]
und umgekehrt jene aus dieser. Der Beweis dieser letztern Congruenz ergiebt sich leicht auf folgende Art. Ist $\mu$ die Anzahl der quadratischen Reste $\alpha_1$, welche zwischen $0$ und $\frac14p$ liegen, so ist (nach VII)
\[
m=2\mu-\tfrac14(p-1)
\]
und die Anzahl der quadratischen Reste $\alpha_2$, welche zwischen $\frac14p$ und $\frac12p$ liegen, ist $=\frac14(p-1)-\mu$. Ist nun $p\equiv1\pmod8$, also die Zahl $2$ quadratischer Rest, so stimmen die Zahlen $2\alpha_1$ und $p-2\alpha_2$ im Complex mit den Zahlen $\alpha_1$ und $\alpha_2$ überein, und bezeichnet man das Product dieser Zahlen mit $A$, so ergiebt sich
\[
2^{\frac{p-1}4}A\equiv(-1)^{\frac14(p-1)-\mu}A\pmod p
\]
und folglich
\[
2^{\frac{p-1}4}\equiv(-1)^\mu\equiv f^{2\mu}\equiv f^{m+\frac14(p-1)}\pmod p
\]
da ferner in diesem Fall $b\equiv0\pmod4$, und folglich
\[
\frac{p-1}4=(a+1)\frac{a-1}4+\frac{bb}4\equiv2\frac{a-1}4\equiv\frac{a-1}2\pmod4
\]
ist, so erhält man die zu beweisende Congruenz
\[
2^{\frac{p-1}4}\equiv f^{m+\frac{a-1}2}\pmod p
\]

Ist dagegen $p\equiv5\pmod8$, also die Zahl $2$ quadratischer Nichtrest, so stimmen die Zahlen $2\alpha_1$ und $p-2\alpha_2$ mit den sämmtlichen zwischen $0$ und $\frac12p$ liegenden quadratischen Nichtresten überein; bezeichnet man ihr Product mit $B$, und das Product der Zahlen $\alpha_1$ und $\alpha_2$ wieder mit $A$, so ist
\[
f\equiv AB,\quad(-1)^{\frac{p-1}4-\mu}2^{\frac{p-1}4}A\equiv B\pmod p
\]
erhebt man diese beiden Congruenzen zum Quadrat, indem man berücksichtigt, dass
\[
ff\equiv-1,\quad2^{\frac{p-1}2}\equiv-1\pmod p
\]
ist, so erhält man
\[
-1\equiv AABB,\quad-AA\equiv BB
\]
und hieraus $A^4\equiv+1$; da nun $A$ ein Product aus quadratischen Resten, also $AA$ ein Product aus biquadratischen Resten und folglich selbst ein biquadratischer Rest ist, so muss $AA\equiv+1$ sein, weil $-1$ ein biquadratischer Nichtrest ist. Hieraus folgt
\SourceLeaf{311}
\[
(-1)^{\frac{p-1}4-\mu}2^{\frac{p-1}4}\equiv AB\equiv f\pmod p
\]
und
\[
2^{\frac{p-1}4}\equiv(-1)^{\mu-1}f\equiv f^{2\mu-1}\equiv f^{m+\frac{p-5}4}\pmod p
\]
da endlich in diesem Fall $b\equiv2\pmod4$, und folglich
\[
\frac{p-5}4=(a+1)\frac{a-1}4+\frac{bb-4}4\equiv2\frac{a-1}4\equiv\frac{a-1}2\pmod4
\]
ist, so erhält man wieder die zu beweisende Congruenz
\[
2^{\frac{p-1}4}\equiv f^{m+\frac{a-1}2}\pmod p
\]
\BookRule{35mm}
